\chapter{Projective geometry}\label{geometry:ch:projective}

\begin{center}\itshape
Metrical geometry is thus a part of descriptive geometry, and descriptive geometry is all geometry.
\upshape(Arthur Cayley, \emph{A Sixth Memoir upon Quantics}, 1859)
\end{center}

Stand between two railway tracks and look along them. The rails are parallel, they never meet, and yet you see them meet, at a point on
the horizon. Every painter of the early Renaissance had to decide what to do about this. The architect Filippo Brunelleschi, in Florence in
the first decades of the fifteenth century, painted the Baptistery on a small panel and made a demonstration of it: according to his
biographer Antonio Manetti, the viewer looked through a hole in the back of the panel at its reflection in a mirror, and the painting
matched the real building seen from the same spot. In 1435 Leon Battista Alberti wrote down the method in \emph{De pictura}: a picture is a
window, and the painter records where the rays from the eye to the scene cross the window. Parallel lines of the floor become lines that run
to a single \emph{vanishing point}, and all the vanishing points lie on one line, the horizon (Figure~\ref{geometry:fig:floor}).

Two centuries later Girard Desargues, an engineer and architect from Lyon, took the painters' idea seriously as mathematics: parallel lines
\emph{do} meet, at a point at infinity, and all the points at infinity form a line. His treatise on conics, the \emph{Brouillon project} of
1639, was printed in some fifty copies and then essentially forgotten. The subject was reborn in a prison camp. Jean-Victor Poncelet, a young
officer in Napoleon's army, was captured in Russia in 1812 and held at Saratov on the Volga until 1814; without books, he rebuilt geometry
from memory and from the idea of projection, and in 1822 published his \emph{Trait\'e des propri\'et\'es projectives des figures}. By 1859
Arthur Cayley could claim, in the words of the epigraph, that projective geometry (``descriptive geometry'' in his vocabulary) contains all of
geometry: distances and angles are projective properties relative to a fixed conic. Appendix~J will make this precise for the hyperbolic plane.

This chapter builds the projective plane from linear algebra, proves the two great incidence theorems of Desargues and Pappus, and then shows
that the projective plane is rigid: over the real numbers, any bijection that sends lines to lines comes from a matrix. The one number that
projections preserve, the cross-ratio, will reappear in every chapter that follows. We end with finite projective planes, where a plane can have
as few as seven points, and with a question about them that was answered only by a long computer search.

\begin{think}
Two triangles $ABC$ and $A'B'C'$ are drawn so that the three lines $AA'$, $BB'$, $CC'$ pass through one point. Extend the sides: $AB$ meets
$A'B'$ somewhere, $BC$ meets $B'C'$, and $CA$ meets $C'A'$. Draw it carefully. What do you notice about these three points? Now imagine the two
triangles are in different planes in space. Can you explain what you saw?
\end{think}

\section{The projective plane}

\begin{definition}\label{geometry:def:projplane}
Let $K$ be a field. The \emph{projective plane} $\mathbb P^2(K)$ is the set of lines through $0$ in $K^3$, that is, of one-dimensional subspaces.
The line spanned by a nonzero vector $p=(x,y,z)$ is written $[p]=[x:y:z]$; these are the \emph{homogeneous coordinates} of the point, defined up
to a nonzero common factor. A \emph{line} of $\mathbb P^2(K)$ is the set of points $[p]$ with $p$ in a fixed plane through $0$ in $K^3$.
\end{definition}

A plane through $0$ has an equation $ax+by+cz=0$ with $(a,b,c)\ne0$, determined up to a factor, so the lines are also labelled by homogeneous
triples $[a:b:c]$. Writing $\ell\cdot p=ax+by+cz$, the point $[p]$ lies on the line $[\ell]$ exactly when $\ell\cdot p=0$.

The points with $z\ne0$ can be written uniquely as $[x:y:1]$: they form a copy of the affine plane $K^2$, the \emph{affine chart}. The
remaining points $[x:y:0]$ form the line $z=0$, the \emph{line at infinity}. An affine line $ax+by+c=0$ becomes the projective line $[a:b:c]$,
which has exactly one point at infinity, $[b:-a:0]$: its direction. Two affine lines are parallel exactly when they have the same direction, that
is, when they meet at infinity. This is the painters' horizon.

\begin{proposition}\label{geometry:prop:incidence}
In $\mathbb P^2(K)$, two distinct points lie on exactly one line, and two distinct lines meet in exactly one point.
\end{proposition}

\begin{proof}
Two distinct points $[p],[q]$ span a two-dimensional subspace $\langle p,q\rangle$, the only plane through $0$ containing both, so the line through
them exists and is unique. Two distinct planes through $0$ in $K^3$ are distinct two-dimensional subspaces, so their sum is $K^3$ and by the
dimension formula their intersection has dimension $2+2-3=1$: it is exactly one point.
\end{proof}

The proof shows the symmetry between points and lines: both are subspaces, points of dimension one and lines of dimension two, and incidence is
inclusion. The cleanest way to make the symmetry exact is the dot product.

\begin{theorem}[Duality]\label{geometry:thm:duality}
The map sending the point $[p]$ to the line $[p]$ (the line with coefficient vector $p$) and the line $[\ell]$ to the point $[\ell]$ is a bijection
between points and lines that preserves incidence: $[p]$ lies on $[\ell]$ if and only if the point $[\ell]$ lies on the line $[p]$. Consequently,
any theorem about points, lines and incidence in $\mathbb P^2(K)$ remains true when the words ``point'' and ``line'', ``lie on'' and ``pass
through'', ``collinear'' and ``concurrent'' are exchanged.
\end{theorem}

\begin{proof}
Both conditions say $\ell\cdot p=0$, and the dot product is symmetric. A theorem about incidences, applied to the images of a configuration
under this bijection, gives the dual statement for the original configuration.
\end{proof}

In coordinates, everything is a cross product. Over any field, define $p\times q$ by the usual formula. Then the line through $[p]$ and $[q]$ is
$[p\times q]$ and the intersection of the lines $[\ell]$ and $[m]$ is $[\ell\times m]$ (Exercise~\ref{geometry:ex:crossproduct}). Three points $[p],[q],[r]$
are collinear exactly when $p,q,r$ lie in one plane, that is, when $\det(p,q,r)=0$.

\begin{figure}[ht]
\centering
\fig{geometry-ch04-floor}
\caption{A tiled floor seen from a height. The picture is the image of the floor under a projective map. The parallel edges of the tiles
meet at the central vanishing point, the diagonals at two others, and all three lie on the horizon, the image of the line at infinity.}
\label{geometry:fig:floor}
\end{figure}

\section{Projective transformations}

An invertible matrix $M\in\mathrm{GL}_3(K)$ maps lines through $0$ to lines through $0$ and planes to planes, so it induces a bijection of
$\mathbb P^2(K)$ that maps lines to lines. Two matrices induce the same map exactly when they are proportional (Proposition~\ref{geometry:prop:frames}
below), so the group of these \emph{projective transformations} is $\mathrm{PGL}_3(K)=\mathrm{GL}_3(K)/K^\times$. The same construction with
$K^2$ gives the projective line $\mathbb P^1(K)=K\cup\{\infty\}$, where $x\in K$ stands for $[x:1]$ and $\infty$ for $[1:0]$, and the group
$\mathrm{PGL}_2(K)$ acting by
\[
x\longmapsto \frac{\alpha x+\beta}{\gamma x+\delta},\qquad \alpha\delta-\beta\gamma\ne0,
\]
with the usual conventions $\infty\mapsto\alpha/\gamma$ and $-\delta/\gamma\mapsto\infty$. These are the M\"obius maps of Appendix~C, over any
field.

The painter's picture is one of these maps. Put the eye at height $h$ above the floor and the picture in a vertical plane at distance $f$ in front
of it. The ray from the eye to the point $(x,y)$ of the floor, where $y$ is the distance beyond the picture plane, crosses the picture at the
fraction $f/(y+f)$ of its length, so it is seen at horizontal position $fx/(y+f)$ and at height $-fh/(y+f)$ relative to the eye. In homogeneous
coordinates this is $[x:y:1]\mapsto[fx:-fh:y+f]$, a linear map with determinant $f^2h\ne0$. The line at infinity of the floor, made of the points
$[x:y:0]$, goes to the points $[fx:0:y]$, which are at height $0$, the height of the eye: the line at infinity of the floor becomes the horizon.

\begin{definition}
Points of $\mathbb P^1(K)$ are in \emph{general position} if they are distinct; points of $\mathbb P^2(K)$ are in general position if no three
of them lie on a line.
A \emph{frame} is an ordered set of $n+2$ points in general position.
\end{definition}

\begin{proposition}[Frames]\label{geometry:prop:frames}
Let $n\in\{1,2\}$. Given two frames of $\mathbb P^n(K)$, there is exactly one projective transformation mapping the first to the second.
In particular a matrix inducing the identity is a scalar.
\end{proposition}

\begin{proof}
It suffices to map the standard frame, made of the points $[e_i]$ for $i\le n+1$ and the point $[e_1+\dots+e_{n+1}]$, to a given frame
$[p_1],\dots,[p_{n+2}]$, and to show that this map is unique; then compose one such map with the inverse of another. The vectors
$p_1,\dots,p_{n+1}$ are linearly independent (for $n=1$, two distinct points; for $n=2$, three points not on a line), so they form a basis. Write
$p_{n+2}=\sum\lambda_ip_i$. Every $\lambda_i$ is nonzero: if $\lambda_j=0$, then $p_{n+2}$ would lie in the span of the $n$ vectors $p_i$ with
$i\ne j$, so the point $[p_{n+2}]$ would coincide with the other point (for $n=1$) or lie on the line through the other two (for $n=2$), against
general position. The matrix $M$ with columns $\lambda_ip_i$ maps
$e_i$ to $\lambda_ip_i$ and $e_1+\dots+e_{n+1}$ to $p_{n+2}$, as required.

For uniqueness, suppose $M$ induces a map fixing the standard frame. Then $Me_i=d_ie_i$ for some scalars $d_i$, so $M$ is diagonal, and
$M(e_1+\dots+e_{n+1})=(d_1,\dots,d_{n+1})$ is proportional to $(1,\dots,1)$: all $d_i$ are equal and $M$ is a scalar matrix.
\end{proof}

So a projective transformation of the plane is determined by the images of four points, and any four points in general position can be moved
to any other four. In particular any quadrilateral can be projected to a square, which is why a photograph of a square tile tells you so
little about the camera.

\section{Desargues' theorem}

Two triangles $ABC$ and $A'B'C'$ are \emph{in perspective from a point} $O$ if the lines $AA'$, $BB'$, $CC'$ pass through $O$, and \emph{in
perspective from a line} if the three points $AB\cap A'B'$, $BC\cap B'C'$, $CA\cap C'A'$ lie on one line.

\begin{theorem}[Desargues]\label{geometry:thm:desargues}
Let $ABC$ and $A'B'C'$ be triangles in $\mathbb P^2(K)$ with $A\ne A'$, $B\ne B'$, $C\ne C'$, and with corresponding sides on distinct lines. If
they are in perspective from a point $O$ different from all six vertices, then they are in perspective from a line.
\end{theorem}

\begin{proof}
Write $O=[o]$, $A=[a]$, $A'=[a']$. The three points $O,A,A'$ are distinct and collinear, so $o=\alpha a+\alpha'a'$ with $\alpha,\alpha'$ both
nonzero (if $\alpha'=0$ then $O=A$). Replacing $a$ by $\alpha a$ and $a'$ by $-\alpha'a'$, which does not change the points, we get $o=a-a'$. In
the same way we choose vectors with
\[
o=a-a'=b-b'=c-c'.
\]
Then $a-b=a'-b'$. This vector is nonzero, since $a=b$ would give $A=B$. It lies in the plane $\langle a,b\rangle$ and in the plane
$\langle a',b'\rangle$, so the point $R=[a-b]$ lies on the lines $AB$ and $A'B'$; since these lines are distinct, $R=AB\cap A'B'$. Likewise
$P=[b-c]=BC\cap B'C'$ and $Q=[c-a]=CA\cap C'A'$. Finally
\[
(a-b)+(b-c)+(c-a)=0,
\]
so the three vectors are linearly dependent and lie in one plane: $P,Q,R$ are collinear.
\end{proof}

The dual statement of Desargues' theorem (Theorem~\ref{geometry:thm:duality}) reads: if two triangles are in perspective from a line, they are in
perspective from a point. So the converse comes for free.

\begin{figure}[ht]
\centering
\fig{geometry-ch04-desargues}
\caption{Desargues' theorem. The triangles $ABC$ and $A'B'C'$ are in perspective from $O$; the intersections $P,Q,R$ of corresponding sides
lie on one line (red).}
\label{geometry:fig:desargues}
\end{figure}

\begin{beauty}
There is a proof without any computation, if you are willing to leave the plane. Suppose the triangles $ABC$ and $A'B'C'$ lie in two different
planes $\Pi$ and $\Pi'$ of space, with $AA'$, $BB'$, $CC'$ meeting at $O$. The lines $AB$ and $A'B'$ lie in the plane $OAB$, so they meet (projectively)
at a point $R$; likewise for $P$ and $Q$. Each of $P,Q,R$ lies on a side of $ABC$, hence in $\Pi$, and on a side of $A'B'C'$, hence in $\Pi'$. Two
distinct planes meet in a line, so $P,Q,R$ are collinear. The plane theorem is the shadow of this one: a flat figure can be lifted to space,
by moving $O$ out of the plane while keeping the other incidences, and then projected back. The vector proof above is this argument in
disguise: the relation $(a-b)+(b-c)+(c-a)=0$ says that $P,Q,R$ lie in the plane spanned by $a-b$ and $b-c$, which plays the role of the second
plane $\Pi'$.
\end{beauty}

\section{Pappus' theorem}

Around the year 320, Pappus of Alexandria recorded in Book VII of his \emph{Collection} a theorem about two lines and six points. It is the
oldest theorem of projective geometry, and it hides an algebraic secret.

\begin{theorem}[Pappus]\label{geometry:thm:pappusline}
Let $\ell$ and $\ell'$ be distinct lines of $\mathbb P^2(K)$ meeting at $O$. Let $A,B,C$ be distinct points of $\ell$ and $A',B',C'$ distinct
points of $\ell'$, all different from $O$. Then the three points
\[
X=AB'\cap A'B,\qquad Y=AC'\cap A'C,\qquad Z=BC'\cap B'C
\]
are collinear.
\end{theorem}

\begin{think}
The proof of Desargues used only that we can add vectors and multiply them by scalars. Does Pappus need more? Keep an eye on every product in
the next proof.
\end{think}

\begin{proof}
Let $o$ span $O$, and choose vectors $u$ and $v$ spanning a point of $\ell$ and a point of $\ell'$, both different from $O$. Then $u,v,o$ is a basis
of $K^3$ (the plane $\langle u,o\rangle$ is $\ell$, which does not contain $[v]$), and we use coordinates in this basis: $\ell$ is the line $y=0$, $\ell'$ is
$x=0$, and $O=[0:0:1]$. A point of $\ell$ other than $O$ has a nonzero first coordinate, so it can be written uniquely as $[1:0:a]$, and
similarly for $\ell'$:
\[
\begin{gathered}
A=[1:0:a],\quad B=[1:0:b],\quad C=[1:0:c],\\
A'=[0:1:a'],\quad B'=[0:1:b'],\quad C'=[0:1:c'],
\end{gathered}
\]
with $a,b,c$ distinct and $a',b',c'$ distinct. The line through $[1:0:a]$ and $[0:1:b']$ is $[(1,0,a)\times(0,1,b')]=[-a:-b':1]$, and in the same way
$A'B=[-b:-a':1]$. Their intersection is $X=[x]$ with
\[
x=(-a,-b',1)\times(-b,-a',1)=\bigl(a'-b',\ a-b,\ aa'-b'b\bigr),
\]
which is nonzero since $a\ne b$. The same computation, with the letters changed, gives
\[
y=\bigl(a'-c',\ a-c,\ aa'-c'c\bigr),\qquad z=\bigl(b'-c',\ b-c,\ bb'-c'c\bigr).
\]
Now compute $x-y+z$. The first two coordinates cancel at once, and the third is
\[
(aa'-b'b)-(aa'-c'c)+(bb'-c'c)=bb'-b'b=0.
\]
So $x,y,z$ are linearly dependent and $X,Y,Z$ are collinear.
\end{proof}

\begin{figure}[ht]
\centering
\fig{geometry-ch04-pappus}
\caption{Pappus' theorem. The three crossing points $X,Y,Z$ of the ``wrong'' pairs of lines are collinear.}
\label{geometry:fig:pappus4}
\end{figure}

\begin{beauty}
Look at the last line of the proof: everything cancels except $bb'-b'b$. Pappus' theorem is true exactly because multiplication is commutative.
The definitions of this chapter make sense over a skew field, such as the quaternions $\mathbb H$, with vectors multiplied by scalars on one
side; Desargues' proof goes through word for word, since it never multiplies two scalars, and the computation above shows that Pappus fails
as soon as $bb'\ne b'b$ (the companion notebook makes this concrete with quaternions). David Hilbert, in his \emph{Grundlagen der
Geometrie} of 1899, turned this around: in an abstract plane satisfying the incidence axioms, Desargues' theorem allows one to build a skew field
of coordinates, and Pappus' theorem holds if and only if that skew field is commutative. A geometric statement about nine points and nine lines
\emph{is} the commutative law.
\end{beauty}

\begin{remark}
Gerhard Hessenberg proved in 1905 that Pappus' theorem implies Desargues' theorem in any projective plane, by a purely synthetic argument.
There are projective planes, satisfying the incidence axioms of Section~\ref{geometry:sec:finiteplanes}, in which Desargues' theorem fails. The simplest is
the Moulton plane (1902): the real plane in which lines of negative slope are bent, their slope doubling where they cross the $y$-axis. See
Hartshorne's \emph{Foundations of Projective Geometry} for both results.
\end{remark}

\section{The cross-ratio}

Projections do not preserve distances, ratios of distances, or even the midpoint of a segment. They do preserve one number attached to four
points on a line.

\begin{definition}\label{geometry:def:crossratio}
Let $a,b,c,d$ be four points of $\mathbb P^1(K)$ with $a,b,c$ distinct. By Proposition~\ref{geometry:prop:frames} there is a unique $g\in\mathrm{PGL}_2(K)$
with $g(a)=\infty$, $g(b)=0$, $g(c)=1$. The \emph{cross-ratio} is $(a,b;c,d)=g(d)\in K\cup\{\infty\}$.
\end{definition}

When the four points are finite, the map is $g(x)=\frac{(c-a)(x-b)}{(c-b)(x-a)}$, a M\"obius map (its determinant is $(c-a)(c-b)(b-a)\ne0$) that
sends $a,b,c$ to $\infty,0,1$. Hence
\begin{equation}\label{geometry:eq:crossratio}
(a,b;c,d)=\frac{(c-a)(d-b)}{(c-b)(d-a)}.
\end{equation}
If one of the points is $\infty$, the formula still holds after cancelling the two factors that contain it; for example $(a,b;c,\infty)=(c-a)/(c-b)$.

\begin{theorem}\label{geometry:thm:crossratio}
Projective transformations of $\mathbb P^1(K)$ preserve the cross-ratio: $(h(a),h(b);h(c),h(d))=(a,b;c,d)$ for $h\in\mathrm{PGL}_2(K)$.
Conversely, if $(a,b;c,d)=(a',b';c',d')$ with $a,b,c$ distinct and $a',b',c'$ distinct, there is a projective map $h$ of $\mathbb P^1(K)$ with $h(a)=a'$,
$h(b)=b'$, $h(c)=c'$, $h(d)=d'$.
\end{theorem}

\begin{proof}
Let $g$ send $a,b,c$ to $\infty,0,1$. Then $g\circ h^{-1}$ sends $h(a),h(b),h(c)$ to $\infty,0,1$, so it is the map used to define the cross-ratio of
the image points, and that cross-ratio is $g(h^{-1}(h(d)))=g(d)$. For the converse, let $g$ and $g'$ be the normalising maps of the two
quadruples; then $h=g'^{-1}\circ g$ sends $a,b,c$ to $a',b',c'$ and $d$ to $g'^{-1}(g(d))=g'^{-1}(g'(d'))=d'$.
\end{proof}

Now we can define the cross-ratio of four points on a line $L$ of the plane. The line $L$ is a plane $V$ through $0$ in $K^3$; choosing a linear
isomorphism $\varphi:K^2\to V$ identifies $\mathbb P^1(K)$ with $L$, and we take the cross-ratio of the corresponding points of $\mathbb P^1(K)$.
Two choices of $\varphi$ differ by an element of $\mathrm{GL}_2(K)$, so by Theorem~\ref{geometry:thm:crossratio} the answer does not depend on the
choice. For the same reason, a projective transformation of the plane preserves the cross-ratio of collinear points.

\begin{proposition}\label{geometry:prop:perspectivity}
Let $L$ and $M$ be lines of $\mathbb P^2(K)$ and $O$ a point on neither. The \emph{central projection} $L\to M$ from $O$, which sends $X$ to
$OX\cap M$, preserves cross-ratios.
\end{proposition}

\begin{proof}
Let $L=[V]$, $M=[W]$ and $O=[o]$. Since $o\notin W$, we have $K^3=W\oplus Ko$; let $\pi:K^3\to W$ be the projection along $o$. For $x\in V$, the
vector $\pi(x)=x-\lambda o$ lies on the line through $[x]$ and $O$ and in $W$, so it represents $OX\cap M$. The restriction $\pi|_V:V\to W$ is linear
and injective, since its kernel is $V\cap Ko=0$ (as $o\notin V$). So the central projection is induced by a linear isomorphism and preserves
cross-ratios.
\end{proof}

This is exactly what the painters need: four equally spaced posts along a road, seen in perspective, are no longer equally spaced, but their
cross-ratio $(0,1;2,3)=\frac{2\cdot2}{1\cdot3}=\frac43$ is the same in the picture as on the road.

Permuting the four points changes the cross-ratio in a controlled way: there are only six values, $\lambda$, $1/\lambda$, $1-\lambda$,
$1/(1-\lambda)$, $(\lambda-1)/\lambda$ and $\lambda/(\lambda-1)$ (Exercise~\ref{geometry:ex:sixvalues}). The most symmetric case deserves a name.

\begin{definition}
Four collinear points form a \emph{harmonic quadruple} if $(A,B;C,D)=-1$; then $D$ is the \emph{harmonic conjugate} of $C$ with respect to $A,B$.
\end{definition}

For $D=\infty$ this means $(c-a)/(c-b)=-1$, that is, $c=(a+b)/2$: the harmonic conjugate of the midpoint is the point at infinity
(Exercise~\ref{geometry:ex:harmonicmidpoint}). The beautiful fact is that harmonic conjugates can be constructed with a straightedge alone. A
\emph{complete quadrangle} is a set of four points $E,F,G,H$ in general position together with the six lines joining them.

\begin{theorem}[Harmonic construction]\label{geometry:thm:harmonic}
Assume $1+1\ne0$ in $K$. Let $EFGH$ be a complete quadrangle, let $A=EF\cap GH$ and $B=EH\cap FG$, and let $C$ and $D$ be the points where the
other two sides $EG$ and $FH$ meet the line $AB$. Then $(A,B;C,D)=-1$. In particular $D$ depends only on $A,B,C$ and not on the quadrangle.
\end{theorem}

\begin{proof}
By Proposition~\ref{geometry:prop:frames} and the invariance of the cross-ratio, we may move $E,F,G,H$ to any four points in general position. Take
$E=[1:1:1]$, $F=[1:-1:1]$, $G=[-1:1:1]$, $H=[-1:-1:1]$, which are in general position because, for instance, $\det(E,F,G)=-4\ne0$ (and similarly for the
other triples). The line $EF$ is $x=z$ and $GH$ is $x=-z$, so $A=[0:1:0]$. The line $EH$ is $y=x$ and $FG$ is $y=-x$, so $B=[0:0:1]$. The line $AB$
is $x=0$. The line $EG$ is $y=z$, meeting it at $C=[0:1:1]$, and $FH$ is $y=-z$, meeting it at $D=[0:1:-1]$. Identify the line $x=0$ with
$\mathbb P^1(K)$ by $[0:s:t]\mapsto[s:t]$. Then $A=\infty$, $B=0$, $C=1$, $D=-1$, and since the normalising map is the identity,
$(A,B;C,D)=-1$. If we start from $A,B,C$ and a different quadrangle, the point $D$ we construct is the unique point with $(A,B;C,D)=-1$
(Theorem~\ref{geometry:thm:crossratio}, applied to the map $g$), so it is the same point.
\end{proof}

When $1+1=0$ the construction collapses: the three diagonal points $A,B$ and $EG\cap FH$ of every quadrangle become collinear, and $C=D$. This is
what happens in the smallest projective plane of all (Exercise~\ref{geometry:ex:fanodiag}).

\section{The fundamental theorem}

A projective transformation maps lines onto lines. Is the converse true? Over $\mathbb C$, complex conjugation of all three coordinates maps lines
onto lines and is not projective (Exercise~\ref{geometry:ex:complexcollineation}). Over the real numbers the answer is yes, and the proof shows how
geometry recovers arithmetic.

\begin{think}
Suppose a bijection of the real plane maps straight lines onto straight lines. With a straightedge and the ability to draw parallels, how would
you construct the point $(x+y,0)$ from $(x,0)$ and $(y,0)$? And the point $(xy,0)$?
\end{think}

\begin{theorem}[Fundamental theorem of projective geometry]\label{geometry:thm:fundamental}
Let $f$ be a bijection of $\mathbb P^2(\mathbb R)$ such that $f(L)$ is a line for every line $L$. Then $f$ is a projective transformation.
\end{theorem}

\begin{proof}
The points of the standard frame $[1:0:0]$, $[0:1:0]$, $[0:0:1]$, $[1:1:1]$ are in general position, and so are their images. Indeed, if
$f(P),f(Q),f(R)$ lie on a line $M$, let $L'=PQ$; then $f(L')$ is a line through $f(P)$ and $f(Q)$, so $f(L')=M\ni f(R)$, and since $f$ is a
bijection, $R\in L'$. By Proposition~\ref{geometry:prop:frames} there is a
projective $g$ agreeing with $f$ on the frame. Replacing $f$ by $g^{-1}\circ f$, we may assume that $f$ fixes the four points, and we must show
that $f$ is the identity.

\emph{Step 1: the affine plane.} The line at infinity $z=0$ contains $[1:0:0]$ and $[0:1:0]$, so its image is the line through these fixed points,
itself. As $f$ is a bijection, it maps the affine plane $\mathbb R^2$ (the points with $z\ne0$) onto itself and maps affine lines onto affine lines.
Two distinct affine lines are parallel exactly when their projective lines meet on the line at infinity, so $f$ maps parallel lines to parallel
lines. In the affine plane, $f$ fixes the origin $[0:0:1]$ and the point $(1,1)=[1:1:1]$, and it fixes the horizontal and vertical directions
$[1:0:0]$ and $[0:1:0]$. So $f$ maps the $x$-axis and the $y$-axis onto themselves, and maps horizontal lines to horizontal lines and vertical lines
to vertical lines. The point $(1,0)$ is the meeting point of the $x$-axis and the vertical line through $(1,1)$, so it is fixed; similarly
$(0,1)$ is fixed.

\emph{Step 2: a function $\sigma$.} Define $\sigma:\mathbb R\to\mathbb R$ by $f(x,0)=(\sigma(x),0)$; it is a bijection because $f$ maps the
$x$-axis bijectively onto itself, and $\sigma(0)=0$, $\sigma(1)=1$. The line through $(1,0)$ and $(0,1)$ is fixed, so its direction is fixed, and
$f$ maps the line through $(y,0)$ of that direction, which meets the $y$-axis at $(0,y)$, to the line through $(\sigma(y),0)$ of the same direction.
Hence $f(0,y)=(0,\sigma(y))$. The point $(x,y)$ is the meeting point of the vertical line through $(x,0)$ and the horizontal line through $(0,y)$,
so $f(x,y)=(\sigma(x),\sigma(y))$.

\emph{Step 3: $\sigma$ is a ring automorphism.} To add, let $P=(0,1)$. The vertical line through $(y,0)$ meets the horizontal line through $P$ at
$(y,1)$. The line through $(y,1)$ parallel to the line from $P$ to $(x,0)$ consists of the points $(y+tx,1-t)$, and meets the $x$-axis at
$(x+y,0)$. Each step (joining two points, drawing a parallel, intersecting) is respected by $f$, and $f$ fixes $P$ and the two axis directions,
so applying $f$ to the construction for $x,y$ gives the construction for $\sigma(x),\sigma(y)$: $\sigma(x+y)=\sigma(x)+\sigma(y)$. To multiply,
the line through $(x,0)$ parallel to the line joining $(1,0)$ and $(0,y)$ consists of the points $(x-t,ty)$ and meets the $y$-axis at $(0,xy)$.
Applying $f$ gives $\sigma(xy)=\sigma(x)\sigma(y)$, using $f(0,y)=(0,\sigma(y))$ from Step~2.

\emph{Step 4: the only automorphism of $\mathbb R$ is the identity.} If $x\ge0$, then $x=s^2$ and $\sigma(x)=\sigma(s)^2\ge0$. So if $x\le x'$ then
$\sigma(x')-\sigma(x)=\sigma(x'-x)\ge0$: $\sigma$ is increasing. Being additive and multiplicative with $\sigma(1)=1$, it fixes every integer and
every rational number. For a real $x$ and rationals $p<x<q$ we get $p\le\sigma(x)\le q$, and letting $p,q\to x$ gives $\sigma(x)=x$.

\emph{Step 5: conclusion.} By Steps~2 and~4, $f$ is the identity on the affine plane. A point at infinity is the direction of some line through
the origin; that line contains two fixed points, so it is mapped onto itself, and its point at infinity is fixed. So $f$ is the identity.
\end{proof}

\begin{remark}
The same proof over any field $K$ shows that a bijection of $\mathbb P^2(K)$ mapping lines onto lines is a projective transformation followed by
a map $[x:y:z]\mapsto[\sigma(x):\sigma(y):\sigma(z)]$ for a field automorphism $\sigma$. Over $\mathbb Q$ and $\mathbb R$ there are no
automorphisms other than the identity; over $\mathbb C$ there is conjugation (and many wild others); over $\mathbb F_{p^k}$ there is the Frobenius
$x\mapsto x^p$. The real case is due to Karl von Staudt (\emph{Geometrie der Lage}, 1847), and the general statement is in Hartshorne's book.
\end{remark}

\section{Finite projective planes}\label{geometry:sec:finiteplanes}

Definition~\ref{geometry:def:projplane} works over any field, including the finite fields $\mathbb F_q$ of Tome~2. Then everything can be counted.

\begin{proposition}\label{geometry:prop:countplane}
$\mathbb P^2(\mathbb F_q)$ has $q^2+q+1$ points and $q^2+q+1$ lines. Every line has $q+1$ points and every point lies on $q+1$ lines.
\end{proposition}

\begin{proof}
There are $q^3-1$ nonzero vectors, and each point is spanned by $q-1$ of them, so there are $(q^3-1)/(q-1)=q^2+q+1$ points. A line is a plane
through $0$, with $q^2-1$ nonzero vectors, so it has $(q^2-1)/(q-1)=q+1$ points. By duality there are as many lines as points and as many lines
through a point as points on a line.
\end{proof}

For $q=2$ we get the \emph{Fano plane}, with $7$ points and $7$ lines (Figure~\ref{geometry:fig:fano}), named after Gino Fano, who used it in 1892 to show
that one of his axioms for projective geometry did not follow from the others. Its points are the seven nonzero vectors of $\mathbb F_2^3$, and
three of them are collinear exactly when they sum to zero.

\begin{figure}[ht]
\centering
\fig{geometry-ch04-fano}
\caption{The Fano plane $\mathbb P^2(\mathbb F_2)$. The seven lines are the six segments and the circle; each line has three points, and three
points form a line when their coordinates add up to $0$ modulo $2$.}
\label{geometry:fig:fano}
\end{figure}

The counting properties suggest an abstract definition. A \emph{projective plane} is a set of points with a family of subsets called lines such
that two distinct points lie on exactly one line, two distinct lines meet in exactly one point, and there exist four points with no three on a
line. If one line of a finite projective plane has $n+1$ points then all lines do, every point is on $n+1$ lines, and there are $n^2+n+1$ points
(Exercise~\ref{geometry:ex:planeorder}); $n$ is the \emph{order} of the plane. The planes $\mathbb P^2(\mathbb F_q)$ have order $q$, a prime power.

\begin{remark}
The planes $\mathbb P^2(\mathbb F_q)$ are not the only finite projective planes: Oswald Veblen and Joseph Wedderburn found non-Desarguesian
planes of order $9$ in 1907. But every known finite projective plane has prime-power order, and whether this must be so is a famous open
problem. Two results are known. Richard Bruck and Herbert Ryser proved in 1949 that if $n\equiv1$ or $2\pmod 4$ and there is a plane of order $n$,
then $n$ is a sum of two squares; this rules out $6$, $14$, $21$, \dots. The first case it does not decide is $n=10$. In 1989 Clement Lam, Larry
Thiel and Stanley Swiercz announced that there is no plane of order $10$, after a computer search that took several years; Lam's account of the
search is listed in the sources. Order $12$ is open.
\end{remark}

Finite planes are not only a curiosity. Their lines form \emph{block designs} used in the design of experiments and in error-correcting codes,
and the seven lines of the Fano plane, written as vectors of length $7$ over $\mathbb F_2$, are exactly the codewords of weight $3$ of the
Hamming code of length~$7$. They also give a party game (Exercise~\ref{geometry:ex:dobble}).

\section{Exercises}

\begin{exercise}[$\star$]\label{geometry:ex:crossproduct}
Show that over any field the line through two distinct points $[p],[q]$ is $[p\times q]$, and the intersection of two distinct lines $[\ell],[m]$
is the point $[\ell\times m]$. Find the intersection of the affine lines $2x+y=3$ and $4x+2y=1$.
\end{exercise}

\begin{exercise}[$\star$]\label{geometry:ex:parallelinf}
Show that the projective lines of two distinct affine lines meet on the line at infinity if and only if the affine lines are parallel. Which
projective transformations map the affine plane onto itself?
\end{exercise}

\begin{exercise}[$\star$]\label{geometry:ex:sixvalues}
Let $\lambda=(a,b;c,d)$ with the four points distinct. Show that exchanging $a\leftrightarrow b$ and $c\leftrightarrow d$ at the same time
does not change the cross-ratio, and that the $24$ orderings of the points give the six values listed in the text. When do some of these six
values coincide?
\end{exercise}

\begin{exercise}[$\star$]\label{geometry:ex:harmonicmidpoint}
Show that $(a,b;c,\infty)=-1$ if and only if $c$ is the midpoint of $a$ and $b$. Use Theorem~\ref{geometry:thm:harmonic} to construct the midpoint of a
segment with a straightedge, given a line parallel to it.
\end{exercise}

\begin{exercise}[$\star$]\label{geometry:ex:dualpappus}
State the dual of Pappus' theorem and draw it.
\end{exercise}

\begin{exercise}[$\star$]\label{geometry:ex:projline}
Show that $\mathbb P^1(\mathbb F_q)$ has $q+1$ points and that $|\mathrm{PGL}_2(\mathbb F_q)|=(q+1)q(q-1)$. Check the answer by computing
$|\mathrm{GL}_2(\mathbb F_q)|$.
\end{exercise}

\begin{exercise}[$\star\star$]\label{geometry:ex:fano168}
Count the frames of the Fano plane and deduce that $\mathrm{PGL}_3(\mathbb F_2)$ has $168$ elements. Check this against the order of
$\mathrm{GL}_3(\mathbb F_2)$, and show that every bijection of the Fano plane mapping lines to lines is in $\mathrm{PGL}_3(\mathbb F_2)$.
\end{exercise}

\begin{exercise}[$\star$]\label{geometry:ex:fanodiag}
Let $E=[1:0:0]$, $F=[0:1:0]$, $G=[0:0:1]$, $H=[1:1:1]$ in $\mathbb P^2(K)$. Show that the three \emph{diagonal points} $EF\cap GH$,
$EG\cap FH$, $EH\cap FG$ are collinear if and only if $1+1=0$ in $K$. Deduce that in the Fano plane the diagonal points of every quadrangle
are collinear, and in $\mathbb P^2(\mathbb R)$ never.
\end{exercise}

\begin{exercise}[$\star\star$]\label{geometry:ex:planeorder}
In a projective plane (abstract axioms of Section~\ref{geometry:sec:finiteplanes}), show that for any two lines $L,M$ there is a point on neither. Deduce
that all lines have the same number $n+1$ of points, that every point lies on $n+1$ lines, and that there are $n^2+n+1$ points.
\end{exercise}

\begin{exercise}[$\star\star$]\label{geometry:ex:dobble}
A card game uses cards printed with symbols such that any two cards share exactly one symbol. Using $\mathbb P^2(\mathbb F_7)$, design a
deck of $57$ cards with $8$ symbols on each card, $57$ symbols in all, every symbol on $8$ cards. Why could you not do the same with $7$
symbols per card and $43$ cards, in the same way?
\end{exercise}

\begin{exercise}[$\star\star$]\label{geometry:ex:complexcollineation}
Show that $[z_0:z_1:z_2]\mapsto[\bar z_0:\bar z_1:\bar z_2]$ maps lines of $\mathbb P^2(\mathbb C)$ onto lines but is not a projective
transformation.
\end{exercise}

\begin{exercise}[$\star\star\star$]\label{geometry:ex:sylvestergallai}
(Sylvester--Gallai) Let $S$ be a finite set of points of the real plane, not all on one line. Show that some line contains exactly two points of
$S$. Is this true in $\mathbb P^2(\mathbb F_q)$, or in $\mathbb P^2(\mathbb C)$?
\end{exercise}

\section{Hints}

\begin{hint}{geometry:ex:crossproduct}
$p\times q$ is orthogonal to $p$ and $q$ for the dot product, over any field (expand the determinant $\det(p,q,r)=(p\times q)\cdot r$).
\end{hint}
\begin{hint}{geometry:ex:parallelinf}
The point at infinity of $ax+by+c=0$ is $[b:-a:0]$. For the second question, the map must preserve the line $z=0$.
\end{hint}
\begin{hint}{geometry:ex:sixvalues}
Use the formula \eqref{geometry:eq:crossratio}. For the swap $b\leftrightarrow c$, expand the numerator of $1-\lambda$.
\end{hint}
\begin{hint}{geometry:ex:harmonicmidpoint}
Cancel the factors containing $d$. For the construction, the line parallel to the segment meets it at infinity.
\end{hint}
\begin{hint}{geometry:ex:dualpappus}
Points become lines; the two lines $\ell,\ell'$ become two points.
\end{hint}
\begin{hint}{geometry:ex:projline}
Proposition~\ref{geometry:prop:frames} with $n=1$: the group acts simply transitively on ordered triples of distinct points.
\end{hint}
\begin{hint}{geometry:ex:fano168}
Choose the four points one by one. For the last question, use the frames proposition and check that a collineation fixing a frame of the Fano
plane fixes every point.
\end{hint}
\begin{hint}{geometry:ex:fanodiag}
Compute the three points with cross products, then one determinant.
\end{hint}
\begin{hint}{geometry:ex:planeorder}
Project $L$ onto $M$ from a point on neither line.
\end{hint}
\begin{hint}{geometry:ex:dobble}
Cards are lines and symbols are points, or the other way round.
\end{hint}
\begin{hint}{geometry:ex:complexcollineation}
The image of the line $[a:b:c]$ is the line $[\bar a:\bar b:\bar c]$. Which points does the map fix?
\end{hint}
\begin{hint}{geometry:ex:sylvestergallai}
Consider all pairs (point $P$ of $S$, line through two or more points of $S$ not containing $P$) and pick one with the smallest positive distance.
\end{hint}

\section{Solutions}

\begin{solution}{geometry:ex:crossproduct}
\idea A line is determined by any nonzero vector orthogonal to two of its points, and the cross product is the standard such vector.
\writeup For all $r$, $(p\times q)\cdot r=\det(p,q,r)$, a polynomial identity valid over any field. Taking $r=p$ and $r=q$ gives $0$, so both
points lie on the line $[p\times q]$. The vector $p\times q$ is nonzero when $p,q$ are linearly independent (some $2\times2$ minor of the
matrix with rows $p,q$ is nonzero), so it defines a line, and it is the unique line through the two points by Proposition~\ref{geometry:prop:incidence}.
The second statement is the dual one, with the same proof. For the example, $(2,1,-3)\times(4,2,-1)=(1\cdot(-1)-(-3)\cdot2,\ (-3)\cdot4-2\cdot(-1),\
2\cdot2-1\cdot4)=(5,-10,0)$. The lines are parallel and meet at the point at infinity $[1:-2:0]$, their common direction.
\end{solution}

\begin{solution}{geometry:ex:parallelinf}
\idea Intersections with the line at infinity are directions.
\writeup The line $[a:b:c]$ meets $z=0$ where $ax+by=0$, at $[b:-a:0]$, which is the direction vector $(b,-a)$ of the affine line. Two distinct
affine lines meet at infinity if and only if they have proportional direction vectors, which means they are parallel. A projective
transformation maps the affine plane onto itself if and only if it maps the line at infinity onto itself, that is, if the matrix maps the plane
$z=0$ to itself: its last row is $(0,0,\ast)$. Scaling the last entry to $1$, these are the maps $v\mapsto Av+b$ with $A\in\mathrm{GL}_2(K)$:
exactly the affine maps of Appendix~B.
\end{solution}

\begin{solution}{geometry:ex:sixvalues}
\idea The formula is a ratio of products; each transposition rearranges the factors.
\writeup Exchanging $a\leftrightarrow b$ in \eqref{geometry:eq:crossratio} gives $\frac{(c-b)(d-a)}{(c-a)(d-b)}=1/\lambda$, and exchanging $c\leftrightarrow d$
also gives $\frac{(d-a)(c-b)}{(d-b)(c-a)}=1/\lambda$; doing both gives $\lambda$ back. Exchanging $a\leftrightarrow c$ and $b\leftrightarrow d$ gives
$(c,d;a,b)=\frac{(a-c)(b-d)}{(a-d)(b-c)}=\frac{(c-a)(d-b)}{(d-a)(c-b)}=\lambda$. The three double transpositions and the identity form a subgroup $V$ of
order $4$ of $S_4$ (the third double transposition is the product of the other two), and $V$ preserves $\lambda$. So the cross-ratio depends
only on the coset of the ordering modulo $V$: at most $24/4=6$ values.
Exchanging $b\leftrightarrow c$: the numerator of $1-\lambda$ is $(c-b)(d-a)-(c-a)(d-b)=(a-b)(d-c)$, so
$1-\lambda=\frac{(a-b)(d-c)}{(c-b)(d-a)}=\frac{(b-a)(d-c)}{(b-c)(d-a)}=(a,c;b,d)$. The transpositions $(ab)$ and $(bc)$ act on $\lambda$ by
$\lambda\mapsto1/\lambda$ and $\lambda\mapsto1-\lambda$; composing them gives the six maps $\lambda$, $1/\lambda$, $1-\lambda$, $1/(1-\lambda)$,
$1-1/\lambda=(\lambda-1)/\lambda$ and $1/(1-1/\lambda)=\lambda/(\lambda-1)$. Values coincide when $\lambda=1/\lambda$, $\lambda=1-\lambda$ or
$\lambda=\lambda/(\lambda-1)$, that is, in the harmonic case $\lambda\in\{-1,2,\frac12\}$ (three values instead of six), and when $\lambda=1/(1-\lambda)$,
that is $\lambda^2-\lambda+1=0$, $\lambda=e^{\pm i\pi/3}$ over $\mathbb C$, where only two values remain. Four points of the complex line with this
cross-ratio are called \emph{equianharmonic}; the vertices of an equilateral triangle and its centre are an example.
\end{solution}

\begin{solution}{geometry:ex:harmonicmidpoint}
\idea With $d=\infty$, the cross-ratio becomes a simple ratio of distances.
\writeup $(a,b;c,\infty)=\frac{c-a}{c-b}=-1$ if and only if $c-a=b-c$, that is $c=\frac{a+b}2$. To construct the midpoint of a segment $AB$
given a line $m$ parallel to it: the line $AB$ meets $m$ at its point at infinity $D$. Choose a point $E$ not on either line, let $F=EA\cap m$ and
$H=EB\cap m$; let $G=AH\cap BF$. These four points form a quadrangle with $EF\cap GH=A$, $EH\cap FG=B$ and $FH=m$, so $FH\cap AB=D=\infty$. By
Theorem~\ref{geometry:thm:harmonic}, $C=EG\cap AB$ satisfies $(A,B;C,\infty)=-1$: it is the midpoint.
\end{solution}

\begin{solution}{geometry:ex:dualpappus}
\idea Dualise word by word.
\writeup Let $L$ and $L'$ be distinct points, and $a,b,c$ distinct lines through $L$, and $a',b',c'$ distinct lines through $L'$, none of them
equal to the line $LL'$. Then the three lines joining $a\cap b'$ to $a'\cap b$, $a\cap c'$ to $a'\cap c$, and $b\cap c'$ to $b'\cap c$ are
concurrent. A drawing with two pencils of three lines shows a hexagon whose three main diagonals meet in a point: this is the degenerate case of
Brianchon's theorem of Appendix~E, just as Pappus is a degenerate case of Pascal's.
\end{solution}

\begin{solution}{geometry:ex:projline}
\idea Count nonzero vectors, then use simple transitivity on ordered triples.
\writeup $\mathbb P^1(\mathbb F_q)$ has $(q^2-1)/(q-1)=q+1$ points. By Proposition~\ref{geometry:prop:frames} with $n=1$, the group acts on ordered triples of
distinct points freely and transitively, so its order is the number of such triples, $(q+1)q(q-1)$. Check:
$|\mathrm{GL}_2(\mathbb F_q)|=(q^2-1)(q^2-q)$ (a first column, then a second column not proportional to it), and dividing by the $q-1$ scalars
gives $(q^2-1)q=(q+1)q(q-1)$.
\end{solution}

\begin{solution}{geometry:ex:fano168}
\idea A frame is chosen point by point, avoiding the lines already drawn.
\writeup The first point can be any of $7$, the second any of the remaining $6$. The third must avoid the line through the first two, which has
$3$ points: $4$ choices. The fourth must avoid the three lines joining the first three. These lines contain the three chosen points and
three more points, one on each side, all distinct, so $6$ points are excluded and $1$ remains. There are $7\cdot6\cdot4\cdot1=168$ frames, and
$\mathrm{PGL}_3(\mathbb F_2)$ acts simply transitively on them by Proposition~\ref{geometry:prop:frames}, so it has $168$ elements. Over $\mathbb F_2$
the only scalar is $1$, so $\mathrm{PGL}_3(\mathbb F_2)=\mathrm{GL}_3(\mathbb F_2)$, of order $(8-1)(8-2)(8-4)=168$. Finally, let $f$ be a
collineation of the Fano plane. Compose it with the projective map that brings $f$ of the standard frame back; the result fixes $[1:0:0]$,
$[0:1:0]$, $[0:0:1]$, $[1:1:1]$. Each of the remaining three points is the third point of a line through two fixed points (for instance
$[1:1:0]$ on the line through $[1:0:0]$ and $[0:1:0]$), so it is fixed too. Hence $f$ is projective.
\end{solution}

\begin{solution}{geometry:ex:fanodiag}
\idea Compute; the answer is the determinant $-2$.
\writeup The line $EF$ is $z=0$ and $GH$ is $x=y$, so $EF\cap GH=[1:1:0]$. The line $EG$ is $y=0$ and $FH$ is $x=z$, so $EG\cap FH=[1:0:1]$.
The line $EH$ is $y=z$ and $FG$ is $x=0$, so $EH\cap FG=[0:1:1]$. These are collinear if and only if
\[
\det\begin{pmatrix}1&1&0\\1&0&1\\0&1&1\end{pmatrix}=-2=0,
\]
that is, if and only if $1+1=0$ in $K$. Every quadrangle can be moved to $EFGH$ by Proposition~\ref{geometry:prop:frames}, which preserves collinearity.
So in $\mathbb P^2(\mathbb F_2)$ the diagonal points of any quadrangle are collinear (in Figure~\ref{geometry:fig:fano}, take the three corners and the
centre: the diagonal points are the three midpoints, on the circle), and in $\mathbb P^2(\mathbb R)$ they never are. This is Fano's axiom of
1892 and its failure.
\end{solution}

\begin{solution}{geometry:ex:planeorder}
\idea A point off both lines gives a bijection between them by projection.
\writeup Let $E,F,G,H$ be four points with no three collinear. If they all lie on $L\cup M$, then exactly two lie on each line (three on one would be
collinear), say $E,F\in L$ and $G,H\in M$, none of them equal to $L\cap M$ (such a point would lie on both lines, giving three on one). The
lines $EG$ and $FH$ are distinct and meet at a point $X$. If $X\in L$, then $X$ is on $L$ and $EG$, so $X=E$ (as $G\notin L$); then $E\in FH$
and $E,F,H$ are collinear, a contradiction. In the same way $X\notin M$. So a point $O$ on neither line always exists. The map $L\to M$,
$P\mapsto OP\cap M$, is well defined (since $OP\ne M$) with inverse $Q\mapsto OQ\cap L$, so $|L|=|M|$; call this number $n+1$. Dually, for a
point $P$ and a line $L$ not through $P$, the lines through $P$ correspond bijectively to the points of $L$, so every point is on $n+1$ lines
(every point lies off some line, since the lines through it are not all the lines). Finally, fix a point $P$. Every other point lies on exactly
one of the $n+1$ lines through $P$, and each of these lines has $n$ points besides $P$, so there are $1+(n+1)n=n^2+n+1$ points.
\end{solution}

\begin{solution}{geometry:ex:dobble}
\idea Two lines of a projective plane share exactly one point: that is the rule of the game.
\writeup Take the points of $\mathbb P^2(\mathbb F_7)$ as symbols and its lines as cards. By Proposition~\ref{geometry:prop:countplane} there are $7^2+7+1=57$
of each, every card (line) has $7+1=8$ symbols, every symbol is on $8$ cards, and two cards share exactly one symbol because two lines meet
in exactly one point. With $7$ symbols per card the same construction would need a projective plane of order $6$, which does not exist by the
Bruck--Ryser theorem quoted in the remark ($6\equiv2\pmod4$ and $6$ is not a sum of two squares). (The theorem forbids only this construction;
smaller decks with the one-common-symbol rule do exist, for instance by removing cards from a larger deck.)
\end{solution}

\begin{solution}{geometry:ex:complexcollineation}
\idea Conjugation respects equations with conjugated coefficients, but it fixes a frame without being the identity.
\writeup If $az_0+bz_1+cz_2=0$ then $\bar a\bar z_0+\bar b\bar z_1+\bar c\bar z_2=0$, so the line $[a:b:c]$ is mapped into the line $[\bar a:\bar b:\bar c]$,
and onto it because conjugation is its own inverse. The map fixes $[1:0:0]$, $[0:1:0]$, $[0:0:1]$ and $[1:1:1]$. If it were projective, it would be
the identity by Proposition~\ref{geometry:prop:frames}. But it sends $[1:i:0]$ to $[1:-i:0]$, a different point.
\end{solution}

\begin{solution}{geometry:ex:sylvestergallai}
\idea L.~M.~Kelly's proof: look at the closest pair of a point and a line not through it. If the line had three points, a closer pair would exist.
\writeup Since the points are not all collinear, there are pairs $(P,L)$ with $P\in S$, $L$ a line through at least two points of $S$, and
$P\notin L$. There are finitely many, so choose one with the smallest distance $\delta>0$ from $P$ to $L$. Let $F$ be the foot of the perpendicular
from $P$ to $L$. Suppose $L$ contained three points of $S$. Two of them, $Q$ and $R$, would lie on the same closed half-line of $L$ starting at
$F$, with $Q$ between $F$ and $R$ (possibly $Q=F$). Let $L'$ be the line $PR$; it contains two points of $S$ and not $Q$. The distance from $Q$ to
$L'$ is at most the distance from $F$ to $L'$, because along the segment $FR$ the distance to $L'$ decreases linearly to $0$ at $R$. The distance from
$F$ to $L'$ is the height of the right triangle $PFR$ from its right angle, $PF\cdot FR/PR<PF=\delta$. So $(Q,L')$ is a closer pair, a contradiction.
Hence $L$ has exactly two points of $S$.

The theorem fails over finite fields: in $\mathbb P^2(\mathbb F_q)$ take $S$ to be all the points; every line has $q+1\ge3$ of them. It also fails over
$\mathbb C$: the nine inflection points of a smooth complex cubic (the Hesse configuration) have the property that the line through any two of them contains a
third. The proof used the order of the real numbers, through distances.
\end{solution}

\section*{Sources and further reading}

\begin{itemize}
\item A.~Cayley, \emph{A sixth memoir upon quantics}, Philosophical Transactions of the Royal Society of London 149 (1859), 61--90, \S230 (the
epigraph, shortened at the end; the full sentence continues ``\dots and reciprocally'').
\item A.~Jones, \emph{Pappus of Alexandria: Book 7 of the Collection}, Springer, 1986 (Pappus' theorem among the lemmas to Euclid's
\emph{Porisms}).
\item J.~V.~Field, \emph{The Invention of Infinity: Mathematics and Art in the Renaissance}, Oxford University Press, 1997 (Brunelleschi's
demonstration as described by Manetti, Alberti's \emph{De pictura} of 1435, and Desargues); M.~Kemp, \emph{The Science of Art}, Yale University
Press, 1990.
\item J.~V.~Field and J.~J.~Gray, \emph{The Geometrical Work of Girard Desargues}, Springer, 1987 (translation of the \emph{Brouillon project} of
1639, and Desargues' theorem as published in A.~Bosse, \emph{Mani\`ere universelle de M.~Desargues pour pratiquer la perspective}, 1648); the
\emph{Dictionary of Scientific Biography} entry on Desargues by R.~Taton (the fifty copies of 1639).
\item J.~J.~Gray, \emph{Worlds Out of Nothing: A Course in the History of Geometry in the 19th Century}, 2nd ed., Springer, 2010 (Poncelet's captivity at
Saratov and the \emph{Trait\'e} of 1822, von Staudt, Cayley's absolute).
\item R.~Hartshorne, \emph{Foundations of Projective Geometry}, Benjamin, 1967 (axiomatic planes, Desargues and Pappus, coordinatising by a skew
field, Hessenberg's theorem, the Moulton plane, the fundamental theorem); D.~Hilbert, \emph{Grundlagen der Geometrie}, Teubner, 1899;
G.~Hessenberg, \emph{Beweis des Desarguesschen Satzes aus dem Pascalschen}, Mathematische Annalen 61 (1905), 161--172; F.~R.~Moulton, \emph{A simple
non-Desarguesian plane geometry}, Transactions of the AMS 3 (1902), 192--195.
\item G.~Fano, \emph{Sui postulati fondamentali della geometria proiettiva}, Giornale di Matematiche 30 (1892), 106--132; O.~Veblen and
J.~H.~M.~Wedderburn, \emph{Non-Desarguesian and non-Pascalian geometries}, Transactions of the AMS 8 (1907), 379--388.
\item R.~H.~Bruck and H.~J.~Ryser, \emph{The nonexistence of certain finite projective planes}, Canadian Journal of Mathematics 1 (1949), 88--93;
C.~W.~H.~Lam, \emph{The search for a finite projective plane of order 10}, American Mathematical Monthly 98 (1991), 305--318.
\item H.~S.~M.~Coxeter, \emph{A problem of collinear points}, American Mathematical Monthly 55 (1948), 26--28 (Kelly's proof of the
Sylvester--Gallai theorem); M.~Aigner and G.~M.~Ziegler, \emph{Proofs from THE BOOK}, Springer, chapter ``Lines in the plane and decompositions of
graphs''.
\item H.~S.~M.~Coxeter, \emph{Projective Geometry}, 2nd ed., Springer, 1987; J.~Richter-Gebert, \emph{Perspectives on Projective Geometry}, Springer,
2011 (a modern book full of pictures and computations with cross products).
\end{itemize}
